\chapter{Tecnologías}

\section{Lenguaje de programación: Python}

Python constituye el lenguaje de programación principal de la solución
propuesta debido a su amplio ecosistema para el desarrollo de
aplicaciones web, procesamiento de datos y aprendizaje automático. En el
proyecto se emplea para implementar la lógica de negocio, el
procesamiento de métricas provenientes de la infraestructura cloud, el
entrenamiento y ejecución del modelo de detección de anomalías basado en
Isolation Forest, así como el desarrollo de los módulos encargados del
análisis de datos, la automatización de procesos y la validación
mediante pruebas unitarias. La elección de este lenguaje responde a su
versatilidad y a la disponibilidad de bibliotecas especializadas que
simplifican la implementación de soluciones de inteligencia artificial.

\section{Framework de desarrollo web: Flask}

El backend de la aplicación fue desarrollado utilizando Flask, un
framework web ligero para Python que permite construir aplicaciones
modulares con una baja sobrecarga de configuración. En esta solución,
Flask se utiliza para implementar la API REST que actúa como
intermediaria entre la interfaz web y los módulos de procesamiento,
facilitando la recepción de solicitudes, la ejecución del análisis de
las infraestructuras cloud y la devolución de los resultados al usuario.
Asimismo, el framework administra las rutas de la aplicación y la
renderización de las vistas utilizadas por el dashboard. La extensión
Flask-CORS complementa esta funcionalidad al habilitar el intercambio
seguro de información entre el frontend y el backend mediante la
configuración de políticas Cross-Origin Resource Sharing (CORS).

\section{Bibliotecas para aprendizaje automático y análisis de datos}

El desarrollo del sistema se apoya en un conjunto de bibliotecas
especializadas para el procesamiento de datos y la implementación del
modelo de inteligencia artificial. La biblioteca Scikit-learn constituye
el núcleo del sistema de detección de anomalías, ya que proporciona la
implementación del algoritmo Isolation Forest utilizado para identificar
patrones de utilización inusuales en los recursos de infraestructura.
Adicionalmente, esta biblioteca incorpora herramientas para la
normalización de características mediante StandardScaler y para la
evaluación del desempeño del modelo durante la etapa experimental.

Por otra parte, NumPy proporciona las estructuras de datos y operaciones
matemáticas necesarias para el procesamiento eficiente de grandes
volúmenes de información, permitiendo realizar cálculos vectorizados,
estadísticas descriptivas y percentiles sobre las métricas obtenidas de
las instancias cloud. Complementariamente, Pandas se utiliza para la
carga, limpieza, transformación y agregación de los datos provenientes
del conjunto Alibaba Cluster Trace 2018, además de facilitar la
generación de las características que posteriormente son utilizadas por
el modelo de aprendizaje automático. Finalmente, Matplotlib permite
representar gráficamente la distribución de las métricas, el
comportamiento de las anomalías detectadas y los resultados obtenidos
durante la evaluación experimental del sistema. Para reducir los tiempos
de ejecución, la persistencia del modelo entrenado y de los objetos de
preprocesamiento se realiza mediante Joblib, evitando la necesidad de
reentrenar el algoritmo en cada ejecución.

\section{Configuración del sistema}

La configuración dinámica del sistema se implementa mediante archivos en
formato YAML procesados a través de la biblioteca PyYAML. Esta solución
permite desacoplar los parámetros de configuración del código fuente,
facilitando la modificación de reglas de recomendación, umbrales de
detección y definiciones de los distintos tipos de instancias sin
necesidad de recompilar o alterar la lógica del sistema. Esta separación
favorece la mantenibilidad y la extensibilidad de la plataforma.

\section{Validación y aseguramiento de la calidad}

Con el objetivo de garantizar el correcto funcionamiento de los
diferentes módulos desarrollados, se emplea el framework Pytest para la
construcción de pruebas unitarias y de integración. Estas pruebas
permiten verificar el comportamiento del motor de recomendaciones, la
clasificación de anomalías y la correcta implementación de las reglas de
negocio. Adicionalmente, la extensión Pytest-cov se utiliza para medir
la cobertura del código fuente, identificando componentes que requieren
una mayor validación durante el proceso de desarrollo.

\section{Integración con Huawei Cloud}

La comunicación con la infraestructura cloud se realiza mediante los SDK
oficiales de Huawei Cloud, los cuales proporcionan acceso programático a
los distintos servicios utilizados por la plataforma. A través de estas
bibliotecas es posible consultar la información de las instancias
Elastic Cloud Server (ECS), recuperar métricas de monitoreo mediante
Cloud Eye, obtener información sobre los volúmenes administrados por
Elastic Volume Service (EVS) y acceder a la configuración de la red
virtual (VPC). Esta integración constituye la fuente principal de
información utilizada durante la etapa de inferencia del modelo,
permitiendo obtener las métricas necesarias para detectar anomalías y
generar recomendaciones de optimización sobre infraestructuras reales.

\section{Tecnologías del frontend}

La interfaz de usuario se desarrolla utilizando tecnologías web estándar
conformadas por HTML5, CSS3 y JavaScript. HTML5 proporciona la
estructura del dashboard y de las diferentes vistas utilizadas por la
aplicación, mientras que CSS3 se emplea para definir el diseño visual,
la adaptación a distintos tamaños de pantalla y la representación
gráfica del estado de las instancias analizadas. Por su parte,
JavaScript implementa la lógica del cliente, permitiendo consumir las
APIs REST desarrolladas en Flask, actualizar dinámicamente la
información mostrada al usuario e incorporar mecanismos de interacción
que facilitan la exploración de los resultados obtenidos durante el
análisis de la infraestructura.

\section{Algoritmo de detección de anomalías}

El componente central del sistema corresponde al algoritmo Isolation
Forest, una técnica de aprendizaje automático no supervisado diseñada
para detectar observaciones anómalas mediante la construcción de árboles
binarios generados aleatoriamente.

El componente de detección de anomalías tiene como objetivo identificar
instancias cuya utilización de recursos presenta un comportamiento
significativamente diferente respecto del patrón de utilización
aprendido a partir de un conjunto de datos de referencia. La finalidad
de esta etapa no consiste directamente en determinar qué recurso debe
ser redimensionado, sino en reducir el universo de instancias que
posteriormente serán evaluadas por el motor de recomendaciones. De esta
manera, el modelo de aprendizaje automático funciona como una primera
etapa de filtrado, mientras que la determinación de una posible acción
de optimización queda a cargo del módulo de recomendaciones.

El modelo utiliza como fuente de entrenamiento el conjunto de datos
Alibaba Cluster Trace 2018, que contiene información de utilización de
recursos obtenida a partir de una infraestructura de producción. A
partir de las métricas disponibles se construye una representación de
cada instancia que posteriormente será utilizada como entrada del
algoritmo.

Las variables consideradas corresponden a la utilización de CPU,
memoria, disco, tráfico de red entrante y tráfico de red saliente.
Debido a que las métricas originales representan observaciones
individuales obtenidas a lo largo del tiempo, se realiza una etapa de
transformación que permite representar el comportamiento de cada
instancia mediante características estadísticas.

Para cada una de las métricas se calculan tres características: media,
percentil 5 y percentil 95. La media permite representar el nivel
habitual de utilización del recurso, mientras que los percentiles
permiten incorporar información sobre los niveles bajo y alto de
utilización sin depender exclusivamente de valores extremos. De esta
forma, cada instancia queda representada mediante un conjunto de
características que describe tanto su utilización habitual como los
niveles de utilización que alcanza en situaciones de menor y mayor
demanda.

Con cinco métricas y tres características derivadas para cada una, la
representación final utilizada por el modelo está compuesta por 15
variables. Esta transformación permite que Isolation Forest trabaje con
una representación homogénea de las instancias y no directamente con las
series temporales originales.

Una vez construidas y normalizadas las características, se procede al
entrenamiento del modelo Isolation Forest. El algoritmo construye un
conjunto de árboles de aislamiento mediante particiones aleatorias del
espacio de características.

El modelo utilizado se configura con 100 árboles, buscando obtener una
representación estable del espacio de características a partir de
múltiples particiones aleatorias. Asimismo, se establece una semilla
aleatoria fija para garantizar que los experimentos puedan ser
reproducidos bajo las mismas condiciones.

Durante el entrenamiento se establece el parámetro contamination,
utilizado por Isolation Forest para determinar la proporción esperada de
observaciones anómalas dentro del conjunto de entrenamiento.

Una vez finalizada la etapa de entrenamiento, el modelo Isolation Forest
queda disponible para analizar nuevas infraestructuras sin necesidad de
volver a entrenarse para cada análisis. Durante la etapa de inferencia,
el sistema obtiene las métricas de utilización correspondientes a las
instancias de la infraestructura que se desea evaluar. Estas métricas
constituyen los datos de entrada sobre los cuales se aplicará el modelo
previamente entrenado.

El proceso de inferencia comienza aplicando a los datos obtenidos el
mismo procedimiento de transformación utilizado durante el
entrenamiento. Para cada instancia se calculan las características
correspondientes a las métricas de CPU, memoria, disco, tráfico de red
entrante y tráfico de red saliente. A partir de cada métrica se obtienen
la media, el percentil 5 y el percentil 95, generando nuevamente las 15
características que representan el comportamiento de la instancia.

Una vez completado el preprocesamiento, las características normalizadas
son introducidas en el modelo Isolation Forest. El algoritmo recorre los
árboles construidos durante la etapa de entrenamiento y determina qué
tan fácilmente puede aislar cada observación respecto del conjunto de
referencia.

Isolation Forest asigna a cada instancia un anomaly score, que
representa el grado en que su comportamiento se diferencia del patrón
aprendido durante el entrenamiento. Este valor permite establecer una
medida continua de anomalía antes de realizar la clasificación final.

En la implementación devuelve el puntaje de cada observación. Los
valores más bajos representan observaciones con un comportamiento más
susceptible de ser considerado anómalo, mientras que valores
relativamente mayores corresponden a observaciones más próximas al
comportamiento considerado normal.

A partir de los puntajes obtenidos, cada instancia es clasificada como
normal o anómala. Para realizar esta clasificación se utiliza el umbral
asociado al modelo, determinado durante su configuración y entrenamiento
en función del parámetro contamination.

Una instancia cuyo puntaje se encuentra por debajo del umbral
establecido es clasificada como anomalía, mientras que aquellas cuyo
puntaje se encuentra por encima del límite son consideradas normales.
Esta separación permite transformar el resultado continuo producido por
Isolation Forest en una decisión binaria que puede ser utilizada por las
etapas posteriores del sistema.

\section{Arquitectura del software}

La aplicación adopta una arquitectura modular que organiza sus
componentes según su responsabilidad funcional. Los módulos de
adquisición y preprocesamiento de datos son responsables de obtener y
preparar las métricas para su análisis, mientras que el módulo de
aprendizaje automático implementa la detección de anomalías mediante
Isolation Forest. De forma complementaria, el motor de recomendaciones
procesa exclusivamente las instancias identificadas como anómalas y
genera propuestas de optimización de recursos. Finalmente, el sistema
incorpora módulos de integración con Huawei Cloud, generación de
reportes y una aplicación web desarrollada en Flask para la interacción
con el usuario. Esta organización modular favorece la mantenibilidad,
escalabilidad y reutilización de los distintos componentes del sistema.
